% =====================================================================
%  ANNOTATION (Russian) — Bachelor's Thesis, Innopolis University 2026
%  Condensed Russian-language companion to the thesis.
%  Times New Roman 14pt, 1.5 spacing, margins 25/20/20/20,
%  continuous page numbering (DOE requirement 11).
% =====================================================================
\documentclass[14pt,a4paper]{extreport}

\usepackage[T2A]{fontenc}
\usepackage[utf8]{inputenc}
\usepackage[russian]{babel}
\usepackage{times}
\usepackage{geometry}
\geometry{a4paper, left=25mm, right=20mm, top=20mm, bottom=20mm}
\usepackage{setspace}
\usepackage{indentfirst}
\usepackage{titlesec}

\onehalfspacing
\setlength{\parindent}{1.25cm}
\pagestyle{plain}   % continuous page numbers, centred in the footer

\titleformat{\chapter}[hang]{\normalfont\large\bfseries}{\thechapter.}{0.5em}{}
\titlespacing*{\chapter}{0pt}{6pt}{12pt}

\begin{document}

% --- title page (placeholder; the official template is used in the docx) ---
\thispagestyle{empty}
\begin{center}
\textbf{Автономная некоммерческая организация высшего образования}\\
\textbf{«Университет Иннополис»}

\vspace{1.5cm}
\textbf{АННОТАЦИЯ}\\
\textbf{НА ВЫПУСКНУЮ КВАЛИФИКАЦИОННУЮ РАБОТУ}\\
\textbf{(БАКАЛАВРСКУЮ РАБОТУ)}\\
\textbf{ПО НАПРАВЛЕНИЮ ПОДГОТОВКИ}\\
\textbf{09.03.01 «Информатика и вычислительная техника»}

\vspace{1cm}
\textbf{НАПРАВЛЕННОСТЬ (ПРОФИЛЬ) ОБРАЗОВАТЕЛЬНОЙ ПРОГРАММЫ}\\
\textbf{«Анализ данных и искусственный интеллект»}

\vspace{1.5cm}
\textbf{Тема:} Использование памяти и графов знаний для оценки\\
действий интеллектуальных агентов

\vspace{1cm}
Выполнил: Пател Кхуш Прадипбхаи\\
Руководитель: проф.\ Бадер Рашид

\vfill
Иннополис, 2026
\end{center}
\clearpage

\setcounter{page}{2}
\tableofcontents
\clearpage

% =====================================================================
\chapter*{Аннотация}
\addcontentsline{toc}{chapter}{Аннотация}

Большие языковые модели всё чаще применяются в роли автоматических судей, оценивающих ответы других моделей. В настоящей работе исследуется, повышает ли качество таких оценок подключение к судье постоянной структурированной памяти о прошлых оцениваниях. Была разработана и реализована система оценивания с памятью MCTS-MAJ, а также протокол оценивания без утечек данных и двухуровневый аудит «замороженной памяти», который доказывает, что память не изменяется во время оценивания. При обычном протоколе память создавала видимость значительного выигрыша; при протоколе без утечек данных, для модели GPT-4o на восьмидесяти отложенных примерах, ни одна из конфигураций не превзошла статистически значимо однопроходного судью без памяти. Аудит дополнительно выявил реальный дефект утечки данных, а исправленный повторный запуск показал, что ранее зафиксированное резкое падение точности при «отравлении» памяти было артефактом измерения. В работе механистически объясняется, почему память не помогла и не навредила на данном наборе данных. Как кажущаяся польза, так и кажущаяся хрупкость памяти в значительной мере являются артефактами протокола оценивания.

\textbf{Ключевые слова:} большая языковая модель в роли судьи, оценивание с памятью, генерация с дополненным поиском, поиск по дереву Монте-Карло, утечка данных при оценивании, аудит замороженной памяти, статистическая значимость.

% =====================================================================
\chapter{Введение}

\section{Актуальность}

Большие языковые модели сегодня используются не только для генерации текста, но и для его \emph{оценивания}. В парадигме «языковая модель в роли судьи» модель получает задание, ответ-кандидат и критерии оценивания и выносит вердикт. Автоматические судьи ранжируют диалоговые системы, оценивают выходные данные программных агентов и формируют сигнал вознаграждения при обучении моделей. Поскольку их вердикты всё сильнее влияют на последующие решения, надёжность автоматического судьи становится первоочередной научной задачей.

Стандартный судья является \emph{однопроходным}: он оценивает каждый пример изолированно и не сохраняет ничего между оцениваниями. Это ограничение мотивирует \emph{оценивание с памятью}, при котором судья подключён к постоянному хранилищу прошлых оцениваний и извлекает релевантные прошлые случаи перед вынесением вердикта. Однако память двойственна: тот же механизм, который поставляет полезный сигнал, может распространять прошлые ошибки судьи, может быть преднамеренно повреждён и, что наиболее тонко, может приводить к утечке информации между связанными примерами оценивания. Поэтому тема работы актуальна: по мере перехода судей с памятью к практическому применению необходим строгий способ измерять, действительно ли память помогает.

\section{Цель и задачи}

Цель работы — спроектировать, реализовать и строго оценить систему оценивания с памятью на основе языковой модели и определить, при каких условиях постоянная память действительно повышает качество оценивания.

Работа построена вокруг исследовательского вопроса: \emph{когда постоянная память помогает судье на основе языковой модели, когда она ухудшает или искажает оценивание, и как можно оценить судью с памятью способом, доказуемо свободным от утечек данных?}

Для ответа поставлены задачи: спроектировать судью с памятью на типизированном графе; реализовать два поисковых компонента на основе поиска по дереву Монте-Карло; разработать протокол оценивания без утечек данных; реализовать двухуровневый аудит замороженной памяти; провести контролируемые эксперименты с разными условиями памяти; определить, сохраняется ли эффект памяти при строгом протоколе, и объяснить полученный результат.

\section{Научная новизна}

Научная новизна работы состоит в следующем. Во-первых, разработан протокол оценивания без утечек данных для судей с памятью, сочетающий разбиение данных по вопросам с замороженным графом памяти. Во-вторых, реализован двухуровневый аудит замороженной памяти: детерминированный снимок состояния «до и после» с криптографическим отпечатком SHA-256 и перехват во время выполнения всех операций записи в память; это первый известный инструмент оценивания, который доказывает, а не предполагает отсутствие утечки. В-третьих, получен корректирующий эмпирический результат: применение аудируемого протокола опровергает ранее правдоподобный вывод о пользе памяти. В-четвёртых, продемонстрировано методологически, что неаудируемый стенд оценивания может сообщать о несуществующих эффектах.

% =====================================================================
\chapter{Краткий обзор литературы}

\emph{Языковая модель в роли судьи.} Парадигма позволяет масштабно оценивать ответы при малой стоимости, однако судьи подвержены систематическим смещениям и стохастичности, и сами нуждаются в тщательной проверке.

\emph{Оценивание с памятью.} Идея продолжает направление генерации с дополненным поиском и работы по агентам с долговременной памятью. В отличие от прежних систем, в данной работе память организована как типизированный граф и изучается именно как возможный источник утечки данных.

\emph{Поиск по дереву Монте-Карло.} Эвристический алгоритм поиска, известный по игровым задачам, в последнее время адаптируется к рассуждениям языковых моделей. В работе он применяется как разведочный поисковый компонент.

\emph{Достоверность оценивания и утечка данных.} Загрязнение и утечка данных — признанные угрозы достоверности оценивания языковых моделей. Память, обновляемая во время оценивания, является прямым каналом такой утечки; методология работы спроектирована так, чтобы закрыть и доказать закрытие этого канала.

% =====================================================================
\chapter{Методология}

Система MCTS-MAJ состоит из трёх компонентов, работающих с общим графом памяти: судьи с памятью, поискового судьи на основе дерева Монте-Карло и модуля поискового извлечения. Память хранится в графовой базе данных Neo4j с пятью типами узлов (Политика, Попытка, Проблема, Исправление, Семантическая категория) и типизированными связями; каждый узел снабжён векторным представлением текста.

Протокол оценивания без утечек данных закрывает два канала утечки. Разбиение по вопросам гарантирует, что ни один вопрос не попадает одновременно в обучающее и тестовое множества. Замораживание памяти во время оценивания гарантирует, что граф не изменяется при прохождении теста.

Аудит замороженной памяти доказывает отсутствие записей в двух независимых слоях. Первый слой берёт детерминированный снимок состояния графа до и после оценивания с отпечатком SHA-256; запуск признаётся свободным от утечек только при совпадении отпечатков. Второй слой во время выполнения перехватывает все методы записи и возбуждает исключение при любой попытке записи. Аудит выявил реальный дефект: путь оценивания записывал оцениваемый тестовый пример обратно в память; дефект был исправлен.

Определены четыре условия памяти: без памяти, самостоятельно сформированная память, эталонная память с истинными метками и «отравленная» память с долей перевёрнутых меток. Каждая точность сопровождается доверительным интервалом Уилсона; сравнения режимов выполняются на парных данных с точным критерием Макнемара и парным бутстрэп-интервалом.

% =====================================================================
\chapter{Реализация}

Система реализована на языке Python. Отдельные модули отвечают за типизированные узлы и встраивания, за управление графом Neo4j и примитивы аудита, за однопроходного судью, за два поисковых компонента и за сборку режимов оценивания. Драйвер тестового стенда выполняет оценивание без утечек данных с включённым аудитом замороженной памяти.

Извлечение решений судьи переведено со свободного текста на структурированный вывод со схемой, что устранило незаметные ошибки разбора. Поскольку оценивание зависит от сети, добавлена обёртка с повторными попытками и экспоненциальной задержкой; пример, который не удаётся обработать после повторов, помечается как ошибочный и исключается из расчёта точности, а не засчитывается как неверный ответ. Это исправление принципиально: именно оно отличает кажущееся резкое падение точности от истинного результата.

Все числа, таблицы и рисунки работы воспроизводимы из опубликованного исходного кода единым сценарием воспроизведения, который проверяет наличие всех артефактов и журналов аудита.

% =====================================================================
\chapter{Оценка и обсуждение}

Эксперименты используют набор данных EvalsBench: 160 примеров по 80 уникальным вопросам, со сбалансированным распределением проходных и непроходных ответов. При протоколе без утечек данных набор разбивается по вопросам на обучающую и тестовую половины по 80 примеров каждая.

При обычном протоколе память создавала видимость значительного выигрыша — порядка шести-десяти процентных пунктов. При аудируемом протоколе без утечек данных, для модели GPT-4o на 80 отложенных примерах, ни одна конфигурация не превзошла статистически значимо однопроходного судью без памяти: точность однопроходного судьи 70,0\,\%, судьи с памятью 65,0\,\%, поискового судьи 70,0\,\%, объединённого режима 71,8\,\%; все доверительные интервалы перекрываются. Эталонная память с истинными метками не превосходит самостоятельно сформированную память; «бедная» память не уступает полной пятиузловой схеме; результат устойчив по трём случайным разбиениям.

В аудируемых экспериментах «отравление» памяти не вызывает значимого снижения точности ни при одной доле повреждения. Это исправляет более ранний неаудируемый результат, согласно которому объединённый режим резко падал до 21,2\,\% при 50\,\% повреждения. При проверке выяснилось, что падение было артефактом: сетевой сбой привёл к ошибкам части примеров, которые засчитывались как неверные ответы. После исправления обработки ошибок результат при 50\,\% повреждения совпадает с эталонным.

\emph{Почему память не помогла и не навредила.} Кажущаяся польза исчезла потому, что значительная её часть была артефактом утечки данных. Память не несёт полезного сигнала, поскольку модель GPT-4o уже работает вблизи предела собственных рассуждений на данной задаче: извлечённое прошлое оценивание не добавляет того, что модель не могла бы получить сама. Это объяснение является гипотезой, согласующейся с данными, а не напрямую проверенным результатом; решающая проверка — более слабая модель-судья, для которой память должна помогать, — оставлена для будущей работы.

% =====================================================================
\chapter{Заключение}

В работе исследовалось, повышает ли постоянная структурированная память качество судьи на основе языковой модели, и как такую систему можно оценить без утечек данных. Спроектирована и реализована система MCTS-MAJ; разработан протокол оценивания без утечек данных и двухуровневый аудит замороженной памяти.

Ответ на исследовательский вопрос состоит в следующем. При обычном протоколе память создавала видимость значительного выигрыша; при аудируемом протоколе без утечек данных этот выигрыш не подтверждается. Память не помогает, когда судья способен решить задачу за счёт собственных рассуждений. Кажущееся резкое падение точности при «отравлении» памяти оказалось артефактом измерения. Воспроизводимый и защищаемый вклад работы является методологическим: разбиение по вопросам, замороженная память и двухуровневый аудит образуют протокол, который доказывает отсутствие утечки данных.

Общий вывод: на данном наборе данных как кажущаяся польза, так и кажущаяся хрупкость памяти в значительной мере являются артефактами протокола оценивания, и аудируемый протокол без утечек данных является необходимым условием любого достоверного утверждения об оценивании с памятью.

Направления дальнейшей работы: существенное расширение набора данных; применение аудируемого протокола к более слабым и более сильным моделям-судьям и к более сложным предметным областям; оформление протокола и аудита как самостоятельного переиспользуемого инструмента оценивания.

% =====================================================================
\chapter*{Список использованной литературы}
\addcontentsline{toc}{chapter}{Список использованной литературы}

Полный список использованной литературы приведён в выпускной квалификационной работе. Основные источники включают работы по парадигме «языковая модель в роли судьи», по генерации с дополненным поиском и памяти языковых агентов, по поиску по дереву Монте-Карло, по загрязнению и утечке данных при оценивании, а также по статистическим методам (доверительный интервал Уилсона, критерий Макнемара, бутстрэп).

\end{document}
